% 19 November 2026. Drafted for Edward Baker with OpenAI Codex assistance.
% E18 helper explicitly dispatched gpt-6.1-sol/ultra; backend unattested.
\section{An obstruction at actual dyadic phases}
\label{sec:actual-phase-recurrence}

The full-circle supremum of Section~\ref{sec:dyadic-group-obstruction}
can be replaced by the common dyadic phase at the actual height.
For the same independently relaxed core rotations, a paid slow-variation
argument still forces failure at some actual height in every length-19
interval within one natural chart. This is an existence statement
about the estimator, without a sign conclusion for the physical field.

All source quantities remain those of
\eqref{af:definitions}--\eqref{af:segment} and
\eqref{dg:groups}--\eqref{dg:matrix}. Thus $q=f-c$ retains the real-profile
value correction, $v$ retains both whole term-error sums and both
entire delta payments, and $D_\kappa$ remains a separate cost. Set
\begin{align}
 C_{\rm act}(x)&=\sum_rR_r(x,u(x)),\qquad u(x)=-\tau(x)\log2,\notag\\
 K_{\rm act}&=1-g-C_{\rm act},&
 \ell_{\rm act}&=\frac{K_{\rm act}}{|E|}-D_\kappa-|c|,\qquad
 m_{\rm act}=\ell_{\rm act}-v .
 \label{ap:allowance}
\end{align}
The radius still takes a supremum over each independent core rotation
$\eta_r$; only the common dyadic phase is now physical.
Reverse triangle in \eqref{dg:identity}, without taking a supremum
over $u$, gives $|EF|\ge K_{\rm act}$. Paying the complete errors proves
\begin{equation}
 \ell_{\rm act}\le|q|,\qquad
 m_E\le m_{\rm dyad}\le m_{\rm act}
       \le m_{\rm dir}=|f-c|-v.
 \label{ap:comparison}
\end{equation}
The divisor is the exact positive $|E|$, even for a negative numerator.
Direct evaluation receives the same coefficients, phases and payments.
Positive complete allowance would require exactly
$C_{\rm act}<1-g-|E|(D_\kappa+|c|+v)$.

\begin{theorem}[Scoped actual-phase recurrence]
\label{thm:actual-phase-recurrence}
On $\mathcal F_0$ of \eqref{af:segment}, every closed interval
$[a,a+19]$ contained in one fixed natural-cutoff chart contains an
actual height $x_*$ with
\begin{equation}
 C_{\rm act}(x_*)-(1-g(x_*))>\frac1{100},\qquad
 K_{\rm act}(x_*)<-\frac1{100},\qquad
 \ell_{\rm act}(x_*),\,m_{\rm act}(x_*)<-\frac1{800}.
 \label{ap:theorem}
\end{equation}
An excluded upper endpoint is interpreted in the old chart by its
continuous one-sided limit; the witness can be chosen strictly before
that endpoint. No failure at every height, at $x_0$ itself, or at a
seam is asserted.
\end{theorem}

\begin{proof}
Retain only the eight prime cores
$\mathcal P_*=\{7,11,13,17,19,23,29,31\}$, and write
$C_*(x,u)=\sum_{p\in\mathcal P_*}R_p(x,u)$.
The finite Fourier-mean argument of the preceding section and
\eqref{af:prime-mass} imply, at frozen height $a$,
\begin{equation}
 \langle C_*(a,\cdot)\rangle>
                  \frac{21}{20}(1-g(a)).
 \label{ap:mean}
\end{equation}
Indeed $\beta_p=b_p$, $a_p=p^{1/5}b_p$,
$h_p\ge(1-g)b_p$, and $\sigma<4/5$.

We first pay for coefficient variation. With $T=x/2$ and
$L_t=L_0+\alpha^2/40$, the explicit carrier derivatives are
\[
 \alpha'=-\frac1{2w^2}-\frac1{(w-1)^2}+\frac1{2w},\qquad
 \alpha''=\frac1{w^3}+\frac2{(w-1)^3}-\frac1{2w^2}.
\]
On $w=v+\ii T$, $2/5\le v\le3/5$ and $7000<T\le15000$,
$1000<|w|/(2\pi)<2500<e^8$ and $|\arg w|<\pi/2$ give
$|\alpha|<6$. Triangle then yields
\[
 |\alpha'|\le\frac1{2T}+\frac3{2T^2}
       <\frac{51}{100T},\qquad
 |\alpha''|<T^{-2},\qquad
 |L_t''|<\frac{51}{100T}+\frac7{20T^2}<\frac3{5T}.
\]
These use only $25/8<\pi<22/7$ and $e>8/3$; for example
$2500\cdot6561<8^8$.
Differentiating $s=s_0+\alpha(s_0)/20$ gives
$s'=-\ii/2-\ii\alpha'(s_0)/40$, hence
\begin{equation}
 |\sigma'|<\frac1{39x},\qquad
 |\tau'+1/2|<\frac1{39x},\qquad
 |g'|<\frac9{125x}<\frac1{10x}.
 \label{ap:velocities}
\end{equation}
For the last bound, conjugation of the fixed carrier logarithm gives
$\log g=\Re[L_t(2/5+\ii T)-L_t(3/5+\ii T)]$.
The horizontal integral of $L_t''$ has length $1/5$, and $dT/dx=1/2$,
so $|(\log g)'|<3/(25x)$; now use $g<3/5$.
The carrier logarithm keeps its fixed analytic branch; the cutoff
floor is not differentiated.

For these eight cores only $d=1,2$ occurs in the convolution.
The original indices at most48 are
\[
 7,14,28;\ 11,22,44;\ 13,26;\ 17,34;\
 19,38;\ 23,46;\ 29;\ 31.
\]
There are16, and each chart retains its $n\le N$ subset.
Each contributes $b_n$ at $m=n$ and $-b_2b_n$ at $m=2n$,
with the additional $n^{1/5}$ for $a_m$.
Using $\log n<4$, $\sigma>3/5$ and $n\ge7$ gives
$b_n n^{-\sigma}<2/5$,
$n^{1/5}b_n n^{-\sigma}<3/5$, and
$b_2 2^{-\sigma}<3/4$. The middle inequality follows from
$(5/3)^{10}<3^5<7^3$.
Consequently, summing over these eight cores only,
\begin{equation}
 B_{\rm abs}:=\sum_m|\beta_m|m^{-\sigma}<\frac{56}{5},
 \qquad
 A_{\rm abs}:=\sum_m|a_m|m^{-\sigma}<\frac{84}{5}.
 \label{ap:mass}
\end{equation}
These triangle bounds are used for variation, not for the physical
field. They include every edge through $m\le2N\le96$, and
$\log m<5$ since $96<(8/3)^5$.

Put $k=(1-g)/(1+g)$ and
$\mathsf Q_p=(1-g)p^{-\sigma}M_p$, so $R_p=\|\mathsf Q_p\|_2$.
Its $m=p2^j$ harmonic is
\[
 m^{-\sigma}
 \operatorname{diag}\bigl(k(\beta_m+ga_m),\beta_m-ga_m\bigr)
 \operatorname{Rot}(ju).
\]
Here $\operatorname{Rot}(v)$ is the real rotation matrix through $v$.
At fixed auxiliary $u$, its $g$-derivative has norm at most
$m^{-\sigma}(2|\beta_m|+|a_m|)$, because
$|k'|\le2$ and $|(kg)'|\le1$.
Its $\sigma$-derivative has norm at most
$\log m\,m^{-\sigma}(|\beta_m|+g|a_m|)$.
Since $\beta_m,a_m$ are constant on a fixed chart, summing these
derivative bounds with \eqref{ap:velocities}--\eqref{ap:mass} gives
\[
 \frac1x\left\{\frac5{39}\frac{532}{25}
                    +\frac1{10}\frac{196}{5}\right\}
 =\frac{6482}{975x}<\frac7x.
\]
Use $|\|A\|_2-\|B\|_2|\le\|A-B\|_2$ and integrate the smooth
matrix derivative, without assuming differentiability of its norm.
Uniformly in $u$ and $a\le x\le a+19$,
\begin{equation}
 |C_*(x,u)-C_*(a,u)|\le\frac{7(x-a)}a,\qquad
 |g(x)-g(a)|\le\frac{x-a}{10a}.
 \label{ap:variation}
\end{equation}
Both are strict for $x>a$ and equality zero at $x=a$.

Finally $\log2=2\int_0^{1/3}(1-r^2)^{-1}\,dr>2/3$ and
\eqref{ap:velocities} give
\[
 u(a+19)-u(a)>
 19\left(\frac12-\frac1{39\cdot14000}\right)\frac23
 =\frac{19}3-\frac{19}{819000}>\frac{44}7>2\pi.
\]
Since $u'>0$, there is a unique $b<a+19$ with
$u(b)=u(a)+2\pi$. The continuous periodic $C_*(a,\cdot)$
attains a maximum, which is at least its strictly bounded mean
\eqref{ap:mean}. Choose $x_*\in[a,b]$ attaining that phase.
With $d=x_*-a<19$, \eqref{ap:variation} proves
\[
 C_*(x_*,u(x_*))-(1-g(x_*))
 >\frac{1-g(a)}{20}-\frac{71d}{10a}
 >\frac1{50}-\frac{1349}{140000}
 =\frac{1451}{140000}>\frac1{100}.
\]
This also covers $d=0$. The full cost is at least $C_*$.
Thus $K_{\rm act}<-1/100$; the exact divisor $0<|E|<8$
gives $K_{\rm act}/|E|<-1/800$, after which the nonnegative
remaining payments prove \eqref{ap:theorem}.
\end{proof}

For the chart endpoints put $l_{33}=x_0$, $l_N=x_N$ for $N\ge34$,
$r_N=x_{N+1}$ for $N\le47$, and $r_{48}=30000$.
Every $r_N-l_N>38$: the first follows from
$(25/8)(4\cdot34^2-1/40)-14001>38$, interior lengths are
$4\pi(2N+1)>38$, and the last exceeds $7248/7$.
Thus the theorem covers the first and last length-19 interval of every
chart. Its witness lies strictly before an excluded old-chart upper
endpoint; the new chart owns the seam. The interval $[x_0,x_0+19]$
already prevents uniform positivity of this estimator on $\mathcal F_0$,
without determining its value at $x_0$. A fixed-reference transfer still
pays the signed jump and whole source payment of \eqref{uc:seams}.

The acquired phase-velocity and coefficient-variation information is
new relative to the auxiliary mean argument. Independent core rotations
remain relaxed, so no joint maximizing core phases are asserted to occur
physically. The exact direct margin's sign and measured evaluation cost
are open. Signed anchor/core aggregation, joint restrictions on the
$\eta_r$, other products, and joint correction or source payments remain
outside this exclusion.

The collision screen of Section~\ref{sec:arithmetic-obstruction}
applies separately to the complete literal source enclosure, the
finite-coefficient/phase-velocity estimates, and a strict complete
margin. A changed source can share coarse velocities or finite
coefficients while lacking the full theta enclosure; an included zero
always defeats the strict margin under an honest enclosure.
The exact quadratic control there has perfect approximation and lacks
theta arithmetic. Conditional and location qualifications of the other
controls remain unchanged. The rest of the two-dimensional finite
strip, effective tail onset and paid overlap, initial verification and
whole intermediate-time barrier remain unacquired. No positive strip
floor, physical zero, new coverage or Newman bound follows.
